\chapter[Exam Background]{Exam Background: Forms, Residues, and Global Obstructions}
\label{ch:exam-background}

The main theorem concerns a torus, a Jacobian, and a cubic.  The proof begins
much earlier, with a lesson from the midterm: a local calculation can fail
globally because a differential form carries an obstruction.
The winding form on \(\R^2\setminus\{0\}\),
\[
  \omega_{\mathrm{wind}}=\frac{-y\,dx+x\,dy}{x^2+y^2},
\]
is closed, but \(\int_{S^1}\omega_{\mathrm{wind}}=2\pi\).  It is locally
\(d\theta\), but no global single-valued angle function exists.

\begin{definition}
An obstruction is a locally invisible quantity that is detected by integration
around a global cycle.  In complex analysis the model obstruction is the residue:
\[
  \frac{1}{2\pi i}\int_\gamma \frac{f'(z)}{f(z)}\,dz.
\]
\end{definition}

The same idea drives the later theory of divisors.  Zeros and poles of a
meromorphic function are local data.  On a compact surface they cannot be chosen
freely, because residues and periods force global balancing laws.

\begin{proposition}[The exam mechanism]
\label{prop:exam-mechanism}
The sequence
\[
  \text{forms} \longrightarrow \text{Stokes} \longrightarrow
  \text{residues} \longrightarrow \text{global constraints}
\]
is the mechanism behind the divisor statements used later.
\end{proposition}

\begin{proof}
Stokes' theorem turns the integral of an exact form over a boundary into an
interior integral.  If a meromorphic form is regular except at finitely many
points, deleting small disks around those points makes Stokes applicable.  Letting
the disk radii shrink leaves only small circle integrals, hence residues.  On a
compact surface the outer boundary is absent, so the sum of all local residues
must vanish.
\end{proof}

\begin{keyidea}{Why this matters}
The statement that principal divisors have degree zero is the divisor version of
the residue theorem applied to \(df/f\).  Abel's theorem is a refined version:
not only the total number of zeros and poles but also their Abel sum is
constrained.
\end{keyidea}

\section*{Exercises}
\addcontentsline{toc}{section}{Exercises}

\exerciseintent{Check, by direct differentiation, that the local winding form is closed away from its singularity.}
\begin{exercise}
Write \(\omega_{\mathrm{wind}}=P\,dx+Q\,dy\).  Compute \(Q_x\), compute
\(P_y\), and subtract them to verify directly that
\[
  d\omega_{\mathrm{wind}}=(Q_x-P_y)\,dx\wedge dy=0
\]
on \(\R^2\setminus\{0\}\).
\end{exercise}

\exerciseintent{Compute the period of the winding form on the basic loop and see the obstruction to exactness.}
\begin{exercise}
Parametrize the unit circle by \(x=\cos t\), \(y=\sin t\).  Compute \(dx\),
\(dy\), the numerator \(-y\,dx+x\,dy\), and then the integral
\(\int_{S^1}\omega_{\mathrm{wind}}\).
\end{exercise}

\exerciseintent{Relate the local order of a meromorphic function to the residue of its logarithmic derivative.}
\begin{exercise}
Let \(f\) be meromorphic near \(p\).  Write \(f=(z-p)^m u(z)\) with
\(u(p)\ne0\).  Compute \(df/f\) explicitly and read off the coefficient of
\((z-p)^{-1}\,dz\), showing that it is \(\ord_p(f)\).
\end{exercise}

\exerciseintent{Turn the compact residue theorem into the degree-zero property of principal divisors.}
\begin{exercise}
Assume \(X\) is compact and \(f\) is a nonzero meromorphic function on \(X\).
Compute
\[
  \sum_p\Res_p(df/f)
\]
in two ways: first by the compact residue theorem, and second using the previous
exercise.  Conclude that \(\degdiv\divisor(f)=0\).
\end{exercise}
